\chapter{Đánh giá công việc và định hướng phát triển}
\textit{Chương này tổng kết mức độ hoàn thành của đồ án, nêu các kết quả nổi bật, chỉ ra những hạn chế còn tồn tại và đề xuất các định hướng phát triển trong giai đoạn tiếp theo.}

\section{Đánh giá kết quả thực hiện}
Dựa trên mục tiêu và phạm vi đã đề ra từ đầu, nhóm đã xây dựng được nền tảng ban đầu cho hệ thống hỗ trợ ôn luyện tiếng Anh ứng dụng AI. Kết quả đạt được cho thấy hệ thống đã hình thành được các thành phần cốt lõi ở cả hai phương diện: hạ tầng ứng dụng web và các mô-đun AI phục vụ học tập.

Ở mức hệ thống, nhóm đã thiết kế được kiến trúc tổng thể, mô hình dữ liệu và các luồng nghiệp vụ chính cho người học, giáo viên hoặc mentor, phụ huynh và quản trị viên. Ở mức chức năng, hệ thống đã hiện thực được các mô-đun quan trọng như quản lý người dùng, tổ chức bài thi thử, theo dõi tiến trình học tập và các thành phần hỗ trợ học tập có tích hợp AI.

\section{Các kết quả nổi bật}
\subsection{Nền tảng ứng dụng web}
Nhóm đã xây dựng được nguyên mẫu ứng dụng web với các nhóm chức năng quan trọng phục vụ bài toán ôn luyện tiếng Anh. Hệ thống hỗ trợ tổ chức dữ liệu bài thi, quản lý thông tin người dùng, lưu lịch sử học tập và tạo nền tảng để tích hợp các dịch vụ AI vào trong quy trình sử dụng thực tế.

Một kết quả đáng chú ý là kiến trúc hệ thống đã được thiết kế theo hướng tách biệt tương đối rõ giữa frontend, backend và các dịch vụ AI. Cách tổ chức này giúp hệ thống thuận lợi hơn trong việc mở rộng, bảo trì và thay thế từng thành phần khi cần thiết.

\subsection{Mô-đun Speaking}
Đối với kỹ năng Speaking, nhóm đã hiện thực pipeline xử lý nhiều giai đoạn
trên nền tảng Lingriser, bao gồm phiên âm lời nói (Whisper-1), phát hiện
khoảng ngừng và ước lượng độ trôi chảy dựa trên word-level timestamps, và
sinh phản hồi ngôn ngữ bằng GPT-4o-mini. Việc kết hợp Whisper và mô hình
ngôn ngữ lớn cho thấy khả năng xây dựng một pipeline đánh giá nói tương đối
toàn diện trong phạm vi đồ án.

\subsection{Mô-đun Writing}
Đối với kỹ năng Writing, nhóm đã thực hiện các bước từ phân tích dữ liệu, chuẩn hóa đầu vào, huấn luyện mô hình Phi-3 bằng phương pháp fine-tuning cho đến triển khai mô hình dưới dạng dịch vụ suy luận. Kết quả bước đầu cho thấy hướng tiếp cận mô hình chuyên biệt cho miền IELTS Writing là khả thi và có tiềm năng hỗ trợ tốt cho bài toán chấm điểm, phân tích bài viết và sinh phản hồi.

Việc triển khai Writing service như một lớp trung gian giữa backend chính và mô hình suy luận cũng là một kết quả có ý nghĩa về mặt kiến trúc. Cách tiếp cận này giúp hệ thống dễ tích hợp, thuận tiện trong giám sát và tạo nền tảng cho các cải tiến tiếp theo.

\section{Lợi ích thực tế của hệ thống}

So với các nền tảng hiện có đã phân tích ở Chương~3, hệ thống mang lại một số lợi ích thực tế rõ rệt:

\begin{enumerate}
    \item \textbf{Tích hợp đánh giá AI cho kỹ năng Nói và Viết trên cùng một nền tảng:} Trong khi phần lớn các nền tảng phổ biến như Study4 không hỗ trợ thi Nói/Viết, hoặc IELTS Online Tests phải dùng giám khảo thật (chi phí cao), hệ thống của nhóm sử dụng pipeline AI hoàn toàn tự động, bao gồm Whisper cho phiên âm, GPT-4o-mini cho phản hồi Speaking và Phi-3 fine-tuned cho chấm điểm Writing, giúp người học nhận phản hồi chi tiết mà không phụ thuộc vào giám khảo.

    \item \textbf{Giảm chi phí ôn luyện:} Các nền tảng như IELTS Online Tests, The Forum và Prep Edu đều yêu cầu trả phí. Hệ thống của nhóm cung cấp miễn phí các tính năng cốt lõi bao gồm làm bài thi thử, chấm điểm AI và phản hồi tự động, giúp mở rộng khả năng tiếp cận cho người học có ngân sách hạn chế.

    \item \textbf{Hỗ trợ đa kỳ thi trên một hệ thống thống nhất:} Thay vì phải sử dụng nhiều nền tảng riêng lẻ cho từng kỳ thi, người học có thể ôn luyện IELTS, TOEIC, TOEFL và SAT trên cùng một giao diện với dữ liệu học tập được tổng hợp, giúp tiết kiệm thời gian chuyển đổi giữa các công cụ và theo dõi tiến độ nhất quán.

    \item \textbf{Hiệu năng đáp ứng yêu cầu thực tế:} Kết quả kiểm thử với k6 cho thấy hệ thống duy trì thời gian phản hồi trung bình dưới 6ms và tỷ lệ thành công trên 98\% ngay cả khi chịu tải 1.000 người dùng đồng thời, xác nhận khả năng phục vụ nhóm người dùng quy mô vừa trong điều kiện vận hành thực tế.

    \item \textbf{Kiến trúc microservices cho phép mở rộng linh hoạt:} Việc tách biệt các service (AI, xác thực, bài thi, thông báo) cho phép nâng cấp hoặc thay thế từng thành phần mà không ảnh hưởng toàn hệ thống. Ví dụ, khi mô hình AI mới ra đời, chỉ cần cập nhật AI service mà không cần sửa đổi backend chính hay frontend.
\end{enumerate}

\begin{table}[H]
\centering
\caption{So sánh lợi thế của hệ thống so với các nền tảng hiện có}
\renewcommand{\arraystretch}{1.4}
\small
\resizebox{\textwidth}{!}{%
\begin{tabular}{|p{0.34\linewidth}|c|c|c|c|c|}
\hline
\textbf{Tiêu chí} & \textbf{Study4} & \textbf{IELTS OT} & \textbf{The Forum} & \textbf{Prep Edu} & \textbf{Hệ thống} \\
\hline
Đánh giá Speaking bằng AI & $\times$ & $\times$ & $\times$ & $\times$ & \checkmark \\
\hline
Chấm Writing tự động & $\times$ & $\times$ & $\times$ & $\times$ & \checkmark \\
\hline
Hỗ trợ $\geq$ 4 kỳ thi & \checkmark & $\times$ & $\times$ & \checkmark & \checkmark \\
\hline
Miễn phí tính năng cốt lõi & \checkmark & $\times$ & $\times$ & $\times$ & \checkmark \\
\hline
Kiến trúc mở rộng được & ? & ? & ? & ? & \checkmark \\
\hline
\end{tabular}
}
\end{table}

\section{Những hạn chế hiện tại}
Mặc dù đã đạt được nhiều kết quả tích cực, hệ thống hiện vẫn còn một số hạn chế.

Thứ nhất, nhiều chức năng mới dừng ở mức nguyên mẫu và chưa được hoàn thiện đến mức sẵn sàng cho môi trường sử dụng thực tế với số lượng người dùng lớn. Một số luồng nghiệp vụ đã được thiết kế nhưng chưa được tích hợp đầy đủ hoặc chưa được tối ưu về trải nghiệm sử dụng.

Thứ hai, các mô-đun AI vẫn phụ thuộc đáng kể vào dịch vụ và API bên thứ ba. Điều này giúp rút ngắn thời gian phát triển trong giai đoạn đầu, nhưng đồng thời cũng làm tăng chi phí vận hành và giảm mức độ chủ động khi mở rộng hệ thống.

Thứ ba, thời gian phản hồi của pipeline Speaking hiện chưa tối ưu do nhiều bước xử lý vẫn diễn ra theo hướng tuần tự. Đây là điểm cần cải thiện nếu hệ thống được phát triển theo hướng phục vụ tương tác gần thời gian thực.

Thứ tư, công tác kiểm thử và đánh giá toàn hệ thống vẫn cần được đầu tư thêm. Dù nhóm đã xác định được định hướng kiểm thử cho cả ứng dụng web và các chức năng AI, độ bao phủ kiểm thử hiện tại vẫn chưa đủ để khẳng định độ ổn định của toàn bộ hệ thống trong nhiều bối cảnh sử dụng khác nhau.


\section{Định hướng phát triển}

\subsection{Hoàn thiện tích hợp hệ thống}
Ưu tiên đầu tiên là hoàn thiện việc tích hợp các mô-đun AI vào ứng dụng web để tạo thành một trải nghiệm sử dụng thống nhất. Điều này bao gồm việc chuẩn hóa giao tiếp giữa backend và các AI service, tối ưu các luồng dữ liệu và bảo đảm kết quả AI được phản ánh rõ ràng, nhất quán trên giao diện người dùng.

\subsection{Tối ưu hiệu năng và khả năng mở rộng}
Hệ thống cần được nâng cấp theo hướng xử lý bất đồng bộ và song song hơn, đặc biệt đối với các tác vụ liên quan đến Speaking và Writing. Việc sử dụng hàng đợi tác vụ, cơ chế cache và các chiến lược tối ưu suy luận sẽ giúp giảm độ trễ, tăng khả năng mở rộng và cải thiện trải nghiệm người dùng.

\subsection{Nâng cao chất lượng các mô-đun AI}
Đối với Speaking, nhóm định hướng cải thiện khả năng phản hồi theo ngữ cảnh, tăng độ tự nhiên của tương tác và mở rộng phân tích theo các tiêu chí đánh giá chi tiết hơn. Đối với Writing, hướng phát triển tiếp theo là mở rộng dữ liệu huấn luyện, thêm các dạng đề thi, chứng chỉ khác như TOEIC, cải thiện chất lượng phản hồi và tăng khả năng giải thích của đầu ra để người học dễ sử dụng kết quả hơn trong quá trình cải thiện kỹ năng.

\subsection{Bổ sung yếu tố con người trong vòng lặp học tập}
Một hướng phát triển quan trọng là kết hợp sâu hơn giữa AI và giáo viên hoặc mentor. Hệ thống có thể mở rộng theo hướng hỗ trợ chuyên gia rà soát kết quả chấm, phản hồi bài viết hoặc bài nói, từ đó tạo ra mô hình học tập kết hợp, trong đó AI đóng vai trò hỗ trợ và con người đóng vai trò thẩm định, định hướng.

\subsection{Mở rộng phạm vi ứng dụng}
Về lâu dài, hệ thống không chỉ dừng lại ở vai trò công cụ ôn luyện cho một số kỳ thi tiếng Anh cụ thể, mà có thể phát triển thành một nền tảng học tập số toàn diện hơn. Trên nền kiến trúc hiện có, hệ thống có thể được mở rộng sang các dạng nội dung học tập mới, các mô hình đánh giá khác và các kịch bản cá nhân hóa sâu hơn theo từng nhóm người dùng.

\section{Kết luận chương}
Đồ án đã đạt được những kết quả có ý nghĩa cả về mặt kỹ thuật lẫn định hướng ứng dụng thực tiễn. Dù hệ thống hiện vẫn còn ở giai đoạn đầu và còn nhiều điểm cần hoàn thiện, các thành phần cốt lõi đã được hình thành tương đối rõ ràng. Đây là cơ sở quan trọng để tiếp tục phát triển hệ thống theo hướng ổn định hơn, thông minh hơn và có giá trị sử dụng cao hơn trong bối cảnh giáo dục số.